
\section{Related Work}
We discuss related work on two key relevant subjects: Visual Art art understanding dataset and Multi-Perspective Art explanation.

\subsection{Visual Art Understanding Datasets}

Existing datasets for computational art understanding are predominantly single-perspective resources, each capturing one interpretive dimension independently of the others.
SemArt~\cite{semart_2018} provides 21,384 WikiArt paintings paired with museum-catalog commentaries that blend iconographic and historical content, enabling cross-modal retrieval but offering no formal, emotional, or narrative perspective.
ArtEmis~\cite{achlioptas2021artemis} collects 439K crowd-sourced affective utterances across 80,031 WikiArt paintings — the largest coverage of emotional response to art — but contains no formal, historical, or scene-level descriptions.
ExpArt~\cite{ExpArt_acl2024} contributes expert formal and historical explanations for approximately 3,500 artworks, while Artpedia~\cite{2019_Artpedia} separates visual sentences from contextual ones for 2,930 paintings.
ArtCaps~\cite{ArtCaps} and earlier captioning datasets~\cite{MM_19_art, Iconographic_IC} target general or iconographic descriptions at smaller scale.
OmniArt~\cite{2017omniart} offers the broadest coverage (432K artworks) but provides only structured metadata rather than natural-language perspectives. For visual question answering on art, Garcia et al.~\cite{garcia_dataset_2020} and ArtQuest~\cite{bleidt_artquest_2024} introduce benchmarks that probe factual and semantic understanding but again operate within a single question-answering register.
KALE~\cite{ijcai2024p848} augments captioning with heterogeneous graph context for richer grounding, but remains a single-caption resource.
The common limitation of these datasets is that each captures one slice of art interpretation: a researcher wishing to study the interplay of formal analysis and affective response must combine incompatible datasets with different scales, annotation conventions, and artwork populations.
MMArt addresses this gap by providing four complementary perspectives — narrative, formal, emotional, and historical — simultaneously for the same 74,234 paintings, with a unified harmonized caption as a fifth field.
Table~\ref{tab:dataset_comparison} summarises the interpretive coverage of existing datasets.

\begin{table}[h]
\centering
\caption{Comparison of art understanding datasets across interpretive dimensions.
\checkmark\,= dedicated annotation; $\sim$\,= partial or metadata only; --\,= absent.
KG = knowledge-graph grounding; Unified = harmonized multi-perspective caption.}
\label{tab:dataset_comparison}
\setlength{\tabcolsep}{4pt}
\begin{tabular}{lrcccccc}
\toprule
\textbf{Dataset} & \textbf{Size} & \textbf{Narr.} & \textbf{Form.} & \textbf{Emot.} & \textbf{Hist.} & \textbf{Unified} \\
\midrule
SemArt~\cite{semart_2018}               & 21K  & \checkmark & --         & --         & \checkmark & -- \\
ArtEmis~\cite{achlioptas2021artemis}    & 80K  & --         & --         & \checkmark & --     & -- \\
Artpedia~\cite{2019_Artpedia}           & 3K   & \checkmark & --         & --         & $\sim$     & -- \\
ExpArt~\cite{ExpArt_acl2024}            & 3.5K & --         & \checkmark & --         & \checkmark & -- \\
ArtCaps~\cite{ArtCaps}                  & 4K   & \checkmark & --         & --         & --        & -- \\
OmniArt~\cite{2017omniart}              & 432K & --         & --         & --         & $\sim$    & -- \\
\midrule
\textbf{MMArt (ours)}                 & \textbf{74K} & \checkmark & \checkmark & \checkmark & \checkmark & \checkmark \\
\bottomrule
\end{tabular}
\end{table}

\subsection{Perspective condition Art explanation}

The interpretive structure of MMArt follows Panofsky's iconological model~\cite{panofsky1955meaning}, which distinguishes pre-iconographic description (subject matter), iconographic analysis (symbolic meaning), and iconological synthesis (cultural significance).
Computationally, Bai et al.~\cite{Bai2021ExplainMT} operationalize multi-topic description by generating separate subject, form, and context paragraphs via a knowledge graph, but remain limited to a few hundred artworks.
GalleryGPT~\cite{MM24GalleryGPT} fine-tunes a vision-language model for formal compositional analysis on the PaintingForm dataset; ArtGPT-4~\cite{Yuan2023ArtGPT4TA} targets general artistic understanding via an LLaVA adapter.
Both confirm that specialized models outperform generalists on individual perspectives, yet neither combines multiple perspectives into a unified dataset.
Knowledge grounding is equally critical for historical interpretation, where purely visual inference leads to frequent factual confabulation~\cite{llm_Hallucination_2023}.
ArtRAG~\cite{ArtRAG_MM2025} demonstrates that retrieval-augmented generation over art-historical knowledge graphs substantially improves factual accuracy and interpretive depth.
MMArt incorporates both insights: each perspective is assigned to its best-suited model --- GalleryGPT for formal analysis, Qwen3-VL~\cite{bai2025qwen3vltechnicalreport} for narrative and historical reasoning,  and ArtEmis grounded utterances for emotional response --- and the historical perspective uses ArtRAG with art context external knowledge ~\cite{ArtRAG_MM2025}, reducing factual hallucination to below 4\% across the dataset.
